\chapter{The Single Invariant}
% OUTLINE Ch1 "The Single Invariant". THE CRUX. Earn-criteria live at hardest:
% #1 the number READ off the earned chain (Vols 1-2), never re-postulated or
% re-derived -- gather, don't re-prove; #2 the read-and-fence a SINGLE
% INSEPARABLE MOVE, per-sentence, BOLD-MATH (as explicit as wanted) /
% FENCED-REFERENT (the device's OWN model number, never the lab's), every
% sentence without exception. The lab-gap stated at its HARDEST exactly where
% the number is read. No Lean identifiers; bracket-not-point; step-not-move.
% Gate: Kodo (earn #1/#2 hardest) + Beastmaster (hardest arrival-read).

The four volumes behind this one all point at a single number, and this
chapter reads it. It does not derive it again --- the deriving is done,
across four books, each step earned at a named site and machine-checked ---
it gathers what was earned and states the result. The reader who wants the
proofs has them in the volumes; what this chapter owes is the number
itself, said as plainly as it can be said, with the one thing that must
travel with it never left behind.

Gather the chain, then, in the order the earlier volumes earned it. At
unit displacement the instrument's own proximity slip reads
\textbf{eighteen} --- a machine-checked equality of the model's counts, not
a measurement of any particle. The target the slip is read against is one
past its floor at the next displacement: four and one, \textbf{five} ---
the same five reached by three separate routes through the model and no
other way. From those two the model's crossing is fixed as the square root
of \textbf{eighteen over five}, an exact quadratic surd built from the
model's own constants and from nothing outside them. The walk of the fourth
volume's geometry closes on that crossing in three earned runs and stops,
leaving two rational walls; carried into the units a coupling is quoted in
--- the model's own reciprocal-at-scale, the same exact arithmetic
throughout --- those walls are \textbf{six hundred forty-eight fifths} and
\textbf{thirteen hundred seventy-seven tenths}: in decimals for the eye,
\textbf{129.6} and \textbf{137.7}. Every number in that sentence is exact,
every one is the model's own, and not one of them is a reading of the
electron the world's instruments weigh.

So the single invariant, stated as what this whole series stands behind:
\[
  129.6 \;<\; \alpha^{-1}_{\text{model}} \;<\; 137.7,
\]
the inverse coupling \emph{of the model this series measures}, bracketed
between two exact ratios of its own counts. The subscript is not decoration.
It is the fence, riding in the symbol itself, because in this volume the
number and its boundary are one inseparable statement: there is no reading
of this bracket that is a reading of the physical electron's coupling, and
the book will not write one.

The reader will want the point inside the walls, and the series' protocol
delivers it under the discipline it has kept since the first volume: the
bracket is what is asserted, and any single value quoted from inside it is
a display for the eye, never a claim. Continued past the earned stop, for
display only, the walls close toward a value near \textbf{137.011} --- and
that number, too, is the model's own. It is not asserted, it is not the
bracket, and above all it is not the laboratory's; it is where the model's
own squeeze appears to be heading, shown because hiding it would be its own
dishonesty, and signed by no one.

Here is the fence at its hardest, at the exact moment the number is read,
because this is where a lesser book would cross. The laboratory measures
the electron's inverse coupling and gets a value near \textbf{137.036}.
The model's bracket contains it; the model's deep display, near 137.011,
sits a breath away from it. That neighborhood --- the model's own number
landing this close to the world's measured one --- is the most remarkable
fact in the entire series, and it is not a claim this book makes. No
argument here bridges the model's value to the laboratory's; none
approaches it asymptotically; none is fitted toward it. The closeness was
found, not sought, and it is left standing as exactly what it is: a
neighborhood, remarkable and unexplained, on the near side of a boundary
this series does not cross. To read the model's bracket as a measurement
of the physical electron's coupling is to take a step no volume of this
series takes --- the step the whole construction was built to make
impossible to take by accident. The number is read; the line beside it is
not.

That the model's number can be pushed nowhere is the reason to trust it,
and the reason is structural, not rhetorical. The eighteen is a counted
slip; the five is its floor and one; the crossing is their surd; the walls
are the walk's own stopping fences; the conversion is the model's exact
reciprocal --- and at no joint in that chain is there a free parameter to
turn, a rounding to choose, a target to steer toward. A number with no
adjustable part cannot be walked to a destination, and the destination was
never named. What lands, lands; and what it landed near is the world's
value, which the model was never told.

The single invariant is now on the page, read off the earned chain,
bracketed in the model's own exact counts, its deep display shown and
unsigned, and the laboratory's value named only to be fenced from it. What
remains is to say what kind of instrument produces a reading like this ---
why it is a \emph{gauge}, and why the word \emph{finite} in this volume's
title is earned and not decorative. That is the next chapter: the gauge
that read the number this chapter stated.
